% !TeX encoding = UTF-8
% Главный файл дипломной работы.
% Перед компиляцией заменить все вхождения {{...}} на свои данные (см. ниже).
\documentclass[a4paper,12pt]{diplom}
\inputencoding{utf8}
%\usepackage{paratype} % Шрифты Paratype — отключены ради совместимости с Overleaf
\DeclareMathSizes{12}{13.4}{11}{10}

\usepackage[left=3cm,right=2cm,top=2cm,bottom=2cm]{geometry}
\usepackage[onehalfspacing]{setspace}
\usepackage{indentfirst}
\setlength{\parindent}{1.25cm}

\usepackage[pdftex]{graphicx}
\graphicspath{{figures/}}
\usepackage{array}
\usepackage{booktabs}
\usepackage{longtable}
\usepackage{makecell}
\usepackage{tikz}
\usepackage[linesnumbered,lined,ruled]{algorithm2e}
\usepackage{listings}
\usepackage{icomma}

\mathtoolsset{showonlyrefs}
\usepackage{microtype}
\usepackage[unicode, pdfborder={0 0 0}, pdfstartview=FitV]{hyperref}

% Часто используемые макросы
\newcommand{\N}{\mathbb{N}}
\newcommand{\Z}{\mathbb{Z}}
\newcommand{\R}{\mathbb{R}}
\renewcommand{\le}{\leqslant}
\renewcommand{\leq}{\leqslant}
\renewcommand{\ge}{\geqslant}
\renewcommand{\geq}{\geqslant}

%%%%%%%%%%%%%%%%%%%%%%%%%%%%%%%%%%%%%%%%%%
%% ДАННЫЕ ТИТУЛЬНОГО ЛИСТА — ЗАМЕНИТЬ {{...}} НА СВОИ
%%%%%%%%%%%%%%%%%%%%%%%%%%%%%%%%%%%%%%%%%%

% Направление подготовки и кафедра — уточнить в учебной части
\Napr{Направление 02.03.02 Фундаментальная информатика и информационные технологии\\Профиль <<Компьютерные науки>>}
\Kafedra{Кафедра теоретической информатики}

% Заведующий кафедрой (статус и ФИО)
\ZavKaf{Заведующий кафедрой,\\ д.\,ф.-м.\,н., профессор}{{{ФИО зав. кафедрой}}}

\DocumentType{\large Выпускная квалификационная работа бакалавра}

% Тема работы
\Title{\Large\bfseries Автоматическая классификация фотографий\\ по визуальному стилю\\ на основе transfer learning}

% Научный руководитель (статус и ФИО)
\Chief{Научный руководитель,\\ {{учёная степень, должность}}}{{{ФИО руководителя}}}

% Автор (группа и ФИО)
\Author{Студент группы {{шифр группы, например ИТ-41БО}}}{{{ФИО автора}}}

\City{Ярославль}
\Year{2026}

\begin{document}

\maketitle

%%%%%%%%%%%%%%%%%%%%%%%%%%%%%%%%%%%%%%%%%%
%% РЕФЕРАТ
%%%%%%%%%%%%%%%%%%%%%%%%%%%%%%%%%%%%%%%%%%
\chapter*{Реферат}

Объём \total{page}~с., \total{chapternum}~гл., \total{fignum}~рис.,
\total{tablenum}~табл., \total{bibnum}~источников, \total{appnum}~прил.

\medskip

\Keywords{компьютерное зрение, классификация изображений, transfer learning, EfficientNet, CLIP, multi-label classification, интерпретируемость, Grad-CAM, embeddings, микросервисная архитектура, Spring Boot, FastAPI, React.}

\medskip

Объект исследования "--- задача автоматической классификации фотографических изображений по визуальному стилю.
Цель работы "--- проектирование и реализация программного комплекса, выполняющего стилевую классификацию загруженных пользователем фотографий на основе подхода transfer learning, и углублённый эмпирический анализ свойств обученной модели.
В~работе спроектирована и реализована микросервисная система из~трёх компонентов (backend на~Spring~Boot, ML"=сервис на~FastAPI и~PyTorch, клиентское веб"=приложение на~React) с~асинхронным взаимодействием через очередь сообщений RabbitMQ и~объектное хранилище MinIO. Собран и~размечен датасет из~4691~фотографии в~восьми стилевых категориях; проведено три итерации дообучения свёрточной нейронной сети EfficientNet"~B0 с~поэтапной коррекцией таксономии классов и~устранением смещения предсказаний; конечная модель v3.1 достигает val~acc=0{,}694 и~macro~F1=0{,}687 на~воспроизводимом разделении.
В~рамках углублённой исследовательской части проведены шесть экспериментов: визуальный анализ внимания модели методом Grad"~CAM, оценка качества обучаемых представлений через t"~SNE/UMAP"=проекции и~метрики разделимости, статистический анализ ручных признаков, калибровка апостериорных вероятностей методом temperature scaling, сравнение с~foundation"=моделью CLIP в~режимах linear probe и~zero"=shot, переход к~multi"=label постановке задачи с~автоматической генерацией дополнительных меток.
Область применения "--- системы курации цифровых фотоколлекций, обучающие материалы по~компьютерному зрению, исследовательские прототипы.

\tableofcontents

%%%%%%%%%%%%%%%%%%%%%%%%%%%%%%%%%%%%%%%%%%
%% ВВЕДЕНИЕ
%%%%%%%%%%%%%%%%%%%%%%%%%%%%%%%%%%%%%%%%%%
\chapternonum{Введение}

\textbf{Актуальность темы.} Объёмы цифровой фотографии растут экспоненциально: профессиональные архивы, коллекции социальных сетей, фотобиблиотеки стоковых агентств содержат миллиарды изображений. Ручная классификация таких массивов по содержательным признакам неосуществима экономически, а~существующие системы автоматической разметки опираются преимущественно на~распознавание объектов и~сцен (<<кошка>>, <<здание>>, <<пляж>>) и~не~описывают эстетических, тональных, композиционных характеристик кадра.

Между тем стилистическая характеристика "--- настроение, цветовая палитра, контраст, освещение "--- является ключевой при~отборе изображений для~дизайн"=проектов, рекламных кампаний, художественных коллекций и~личных архивов. Развитие методов глубокого обучения и~накопление общедоступных предобученных моделей сделало решение задачи стилевой классификации технически достижимым в~условиях ограниченного датасета.

Подход, основанный на~transfer learning, при~котором переиспользуется уже обученная на~миллионах изображений ImageNet свёрточная сеть, а~дообучается лишь финальная классифицирующая часть, является оптимальным по~балансу качества и~вычислительной стоимости и~рассматривается в~настоящей работе.

\textbf{Цель работы.} Разработать программный комплекс, позволяющий пользователю загружать фотографии и~получать их~автоматическую стилевую классификацию с~оценкой уверенности модели, осуществлять поиск похожих изображений в~собственной коллекции и~фильтрацию по~стилю; провести углублённый экспериментальный анализ обученной модели с~акцентом на~интерпретируемость, качество представлений и~устойчивость метрик.

\textbf{Задачи работы.}
\begin{enumerate}[label=\arabic{enumi})]
\item провести обзор современных методов классификации изображений, существующих программных продуктов"=аналогов и~технологического стека; обосновать выбор архитектуры и~инструментальных средств;
\item собрать и~разметить датасет фотографий по~стилевым категориям; провести предобработку, балансировку и~дедупликацию;
\item обучить и~оценить нейросетевую модель классификации; провести несколько итераций пересмотра таксономии классов по~результатам диагностики;
\item спроектировать многокомпонентную программную архитектуру с~асинхронной обработкой запросов на~основе очереди сообщений и~объектного хранилища;
\item реализовать REST"=API для~регистрации пользователей, загрузки изображений, получения результатов классификации, поиска по~стилю и~поиска похожих изображений;
\item реализовать клиентское веб"=приложение;
\item обеспечить контейнеризацию и~автоматизированное развёртывание системы;
\item провести шесть углублённых исследовательских экспериментов: визуализацию внимания модели, анализ embedding"=пространства, статистический анализ ручных признаков, калибровку вероятностей, сравнение с~foundation"=моделью CLIP, расширение задачи до~multi"=label;
\item сформулировать направления дальнейшего развития на~основе полученных результатов.
\end{enumerate}

\textbf{Объект исследования} "--- автоматическая стилевая классификация цифровых фотографий. \textbf{Предмет исследования} "--- архитектуры свёрточных нейронных сетей применительно к~задаче с~ограниченным обучающим набором и~субъективной таксономией классов, методы интерпретации их~поведения, способы согласования постановки задачи с~природой стилевых атрибутов.

\textbf{Методы исследования.} Подход transfer learning на~основе свёрточных нейронных сетей с~предобучением на~ImageNet; стратифицированное разделение данных с~детерминированным seed; двухфазное дообучение модели; weighted random sampling для~балансировки классов; cross"=entropy и~binary cross"=entropy loss; cosine annealing learning rate; визуализация активаций методом Grad"~CAM; снижение размерности эмбеддингов методами t"~SNE и~UMAP; статистический анализ распределений (ANOVA, тест Краскела"=Уоллиса); калибровка вероятностей методом Guo et~al.\,(2017); линейный probe и~zero"=shot классификация на~основе модели CLIP.

\textbf{Научная и~практическая новизна.} Научная составляющая "--- адаптация подхода transfer learning к~задаче стилевой классификации с~субъективно определёнными классами; формальное обоснование операциональной определимости отдельных классов через статистики низкоуровневых признаков; экспериментальная демонстрация устойчивости качественных выводов (иерархия сложности классов, структура ошибок) к~случайному разбиению данных. Практическая "--- завершённое программное решение с~открытым исходным кодом, готовое к~демонстрации и~расширению.

\textbf{Структура работы.} Работа состоит из~введения, шести глав, заключения, списка использованных источников и~приложений. В~первой главе проводится постановка задачи, обзор методов, выбор архитектуры, обзор существующих программных продуктов и~технологического стека. Во~второй главе описывается архитектура спроектированной системы. В~третьей главе раскрываются ключевые решения, принятые при~реализации. В~четвёртой главе изложена методология сбора и~предобработки датасета, процедура дообучения модели и~результаты трёх итераций обучения. В~пятой главе представлены результаты функционального тестирования и~углублённого исследовательского анализа модели (шесть экспериментов). В~шестой главе обсуждается развёртывание системы, информационная безопасность и~направления дальнейшего развития. В~заключении сводятся итоги, приводятся достигнутые показатели, перечисляются возможные направления развития.

%%%%%%%%%%%%%%%%%%%%%%%%%%%%%%%%%%%%%%%%%%
%% ГЛАВА 1
%%%%%%%%%%%%%%%%%%%%%%%%%%%%%%%%%%%%%%%%%%
\chapter{Аналитическая часть}

\section{Постановка задачи классификации}
\label{sec:problem}

Формально задача рассматривается в~двух постановках "--- классической single"=label (одной фотографии соответствует один стилевой ярлык) и~расширенной multi"=label (одной фотографии может соответствовать несколько ярлыков с~независимыми степенями уверенности). На~ранних итерациях работы используется single"=label постановка как более простая и~поддерживаемая стандартной средой обучения; в~ходе анализа результатов (см.~гл.~\ref{ch:results}) проводится обоснованный переход к~multi"=label.

В~single"=label постановке выходом модели является вектор $\hat{y} \in \R^C$ постмягкомаксных вероятностей, где $C$~--- количество классов; предсказанием считается $\arg\max_c \hat{y}_c$. В~multi"=label постановке выходом является вектор сигмоид"=активаций; финальное решение для~каждого класса принимается на~основе предобученного порога $\theta_c$.

Целевыми классами выбраны восемь визуальных стилей, операционально определимых через статистики низкоуровневых признаков изображения (см.~табл.~\ref{tab:classes}):

\begin{table}[!ht]
\caption{Целевые стилевые классы и их операциональные характеристики}
\label{tab:classes}
\centering
\begin{tabular}{lp{9.5cm}}
\toprule
\textbf{Класс} & \textbf{Визуальные характеристики} \\
\midrule
\texttt{airy}        & Высокая яркость, мягкий рассеянный свет, светлая палитра, низкий контраст \\
\texttt{dark}        & Низкая средняя яркость, глубокие тени, преобладание тёмных тонов \\
\texttt{dramatic}    & Высокий контраст, выраженная светотень, эмоционально насыщенная композиция \\
\texttt{golden\_hour}& Тёплая цветовая температура, мягкий направленный свет, оранжево"=жёлтые тона \\
\texttt{minimalist}  & Малое количество объектов, большая доля негативного пространства, низкий контраст \\
\texttt{monochrome}  & Около нулевая насыщенность (чёрно"=белая или сильно десатурированная палитра) \\
\texttt{neon}        & Высокая насыщенность ярких источников света на~тёмном фоне (киберпанк, синтвейв) \\
\texttt{vintage}     & Приглушённые тона с~тёплым оттенком, эффект плёнки, пониженная резкость \\
\bottomrule
\end{tabular}
\end{table}

Эта таксономия "--- результат двух итераций пересмотра (см.~гл.~\ref{ch:dataset}). Первоначальный набор включал классы \texttt{moody} и~\texttt{street}, которые были исключены как методологически неудовлетворительные: первый систематически пересекался с~тремя другими (\texttt{dark}, \texttt{dramatic}, \texttt{vintage}) и~не~имел операциональной формулировки, второй являлся жанровой, а~не~стилевой категорией. Введённые взамен классы \texttt{monochrome} и~\texttt{neon} обладают строгим операциональным определением через насыщенность.

Все классы определены так, чтобы быть приближённо ортогональными в~пространстве низкоуровневых признаков, что позже подтверждается статистическим анализом (см.~раздел~\ref{sec:scorecard}). Тем не~менее, на~уровне отдельных изображений возможны совмещения "--- например, кадр может быть одновременно \texttt{monochrome} и~\texttt{minimalist}, "--- что мотивирует переход к~multi"=label постановке в~разделе~\ref{sec:multilabel}.

\section{Обзор методов глубокого обучения}
\label{sec:methods}

\subsection{Обучение модели <<с нуля>>}

Прямое обучение глубокой свёрточной сети, начиная со~случайной инициализации весов, требует датасета порядка $10^5$"=$10^6$ размеченных изображений на~класс и~значительных вычислительных ресурсов. Для~задачи с~восемью классами и~ручной разметкой такой подход неприменим: ни~по~объёму данных, ни~по~вычислительному бюджету. Кроме того, подобное обучение не~использует уже накопленные представления, что является расточительным.

\subsection{Feature extraction}

Промежуточный подход, при~котором выходы предпоследнего слоя предобученной сети интерпретируются как векторные признаки, поверх которых обучается простой классификатор (логистическая регрессия, SVM, k"=ближайших соседей). Метод вычислительно дёшев, но~ограничен по~качеству: признаки ImageNet ориентированы на~распознавание категорий объектов, а~не~стилевых атрибутов. Адаптация представлений к~новому распределению задач невозможна "--- веса замораживаются полностью.

В~настоящей работе этот подход реализован в~двух разновидностях как один из~элементов сравнительного анализа: feature extraction поверх ImageNet"=весов EfficientNet"~B0 (см.~раздел~\ref{sec:scorecard}) и~поверх замороженного CLIP ViT"~B/32 (см.~раздел~\ref{sec:clip}).

\subsection{Fine-tuning (выбранный подход)}

Замена классифицирующего слоя предобученной сети, размораживание последних свёрточных блоков и~дообучение на~целевом датасете с~пониженным learning rate. Подход обеспечивает адаптацию представлений к~стилевым признакам при~сохранении общей иерархии признаков, выученной на~миллионах изображений ImageNet. Для~датасетов среднего размера ($10^3$"=$10^4$ изображений) он~является оптимальным компромиссом между качеством и~стоимостью обучения.

В~настоящей работе применена двухфазная схема fine"=tuning'а: на~первой фазе backbone полностью замораживается, обучается только классифицирующая голова при~повышенном learning rate, что~стабилизирует функционал потерь; на~второй фазе все слои размораживаются, learning rate понижается, применяется cosine annealing"=расписание. Подробное описание процедуры приведено в~разделе~\ref{sec:training}.

\subsection{Дополнительные техники}

Помимо fine"=tuning'а, в~работе применены следующие техники, направленные на~стабилизацию обучения и~повышение качества:
\begin{itemize}
\item \textbf{WeightedRandomSampler} "--- балансирующая выборка минибатчей с~весами, обратно пропорциональными частотам классов; нейтрализует bias предсказаний при~остаточном дисбалансе.
\item \textbf{Label smoothing} с~параметром $\varepsilon=0{,}1$ "--- сглаживание истинной метки, препятствует переуверенности модели и~улучшает калибровку.
\item \textbf{Mixed precision training} (\texttt{torch.amp.autocast}) "--- ускорение обучения на~GPU Tesla T4 в~1{,}5"~2 раза при~сохранении численной устойчивости.
\item \textbf{Early stopping} с~патиенсом семь эпох по~val~accuracy "--- предотвращает переобучение и~экономит вычислительное время.
\item \textbf{Безопасные для~стиля аугментации}: RandomResizedCrop, HorizontalFlip, ограниченный ColorJitter (без~RandomGrayscale, без~сильного hue jitter). Аугментации не~должны уничтожать классифицирующий сигнал (например, RandomGrayscale превратил~бы любую фотографию в~\texttt{monochrome}, нарушив разметку).
\end{itemize}

\section{Выбор архитектуры нейронной сети}
\label{sec:arch}

В~качестве кандидатов рассмотрены три семейства свёрточных сетей с~предобучением на~ImageNet, относящиеся к~разным архитектурным эпохам и~целевым нишам. Их~сводные характеристики приведены в~таблице~\ref{tab:arch}.

\begin{table}[!ht]
\caption{Сравнение архитектур"=кандидатов для backbone классификатора}
\label{tab:arch}
\centering
\begin{tabular}{lrrrr}
\toprule
\textbf{Архитектура} & \textbf{Парам., M} & \textbf{Top"~1, ImageNet} & \textbf{Размер чекп., МБ} & \textbf{Инф., CPU, мс} \\
\midrule
ResNet"~50            & 25{,}6 & 76{,}1\,\% & 100 & $\sim$\,150 \\
MobileNet"~V3 Large   & 5{,}4  & 75{,}2\,\% & 22  & $\sim$\,35  \\
EfficientNet"~B0      & 5{,}3  & 77{,}7\,\% & 16  & $\sim$\,50  \\
\bottomrule
\end{tabular}
\end{table}

В~качестве основного backbone выбрана \textbf{EfficientNet"~B0}~\cite{Tan:2019} по~трём причинам.

\paragraph{Компромисс точности и размера.} При~количестве параметров, сравнимом с~MobileNet"~V3 Large, EfficientNet"~B0 обеспечивает на~$\sim$\,2{,}5\,пп более высокую точность Top"~1 на~ImageNet и~превышает ResNet"~50 на~$\sim$\,1{,}6\,пп при~меньшем в~пять раз количестве параметров. Эффект достигается за~счёт \emph{compound scaling} "--- согласованного масштабирования глубины сети, ширины слоёв и~входного разрешения.

\paragraph{Эффективность инференса на CPU.} Размер чекпоинта 16~МБ и~время инференса $\sim$\,50~мс на~процессоре общего назначения позволяют выполнять предсказания в~production без~использования GPU, что существенно упрощает развёртывание и~сокращает эксплуатационные затраты.

\paragraph{Зрелая реализация в torchvision.} Модель доступна в~библиотеке \texttt{torchvision.models} как \verb|efficientnet_b0(weights=EfficientNet_B0_Weights.IMAGENET1K_V1)|. Один вызов возвращает полностью инициализированный экземпляр сети, что упрощает экспериментальный цикл.

Для контрольного сравнения дополнительно проведено обучение \textbf{ResNet"~50}; результаты см.\ в~разделе~\ref{sec:training}. Сравнение с~архитектурой принципиально иной парадигмы "--- visual transformer "--- проведено в~рамках раздела~\ref{sec:clip} (CLIP ViT"~B/32).

\section{Обзор существующих программных продуктов}
\label{sec:competitors}

Задача автоматического тегирования фотографических изображений по~стилю не~освоена коммерческими продуктами в~явном виде. Существующие решения покрывают смежные ниши, но~отличаются по~ключевым характеристикам "--- степени автоматизации, открытости модели, ориентированности на~стиль (а~не~на~содержимое), возможности self"=hosting.

\paragraph{Pinterest, Are.na, Cosmos.} Платформы визуального discovery и~ручной курации коллекций. Pinterest использует автоматическое распознавание объектов через внутреннюю модель Pinterest Lens (одежда, мебель, еда), но~таксономия стилей в~явном виде отсутствует; теги генерируются пользователями. Are.na и~Cosmos автоматической классификацией не~занимаются вовсе "--- ценность платформ в~человеческом отборе.

\paragraph{Google Photos, Apple Photos.} Встроенные в~экосистемы платформы автоматической классификации фотографий пользователя. Распознают \emph{содержание} (объект, сцена), не~предоставляют API для~извлечения признаков или~таксономии стилей, не~допускают self"=hosting.

\paragraph{Adobe Lightroom AI tagging (Adobe Sensei).} Профессиональный продукт с~автоматическим тегированием. Теги преимущественно технические (\emph{Low light}, \emph{Sharp}, \emph{High contrast}) или~предметные (\emph{Mountain}, \emph{Indoor}). Стилевые понятия (\emph{vintage}, \emph{moody}) присутствуют, но~их~надёжность не~декларируется метриками "--- это маркетинговая функция платной подписки. Модели и~внутренняя реализация закрыты.

\paragraph{Системы style transfer (Prisma, DeepArt, Lensa AI).} Это generative"=системы, решающие \emph{обратную} задачу: применение стилевых характеристик одного изображения к~содержимому другого. Концептуально полезны "--- их~работоспособность подтверждает, что~стиль есть вычислимая характеристика, отделимая от~содержимого, но~задача классификации в~них не~ставится.

\paragraph{Академические работы.} Стилевая классификация изображений "--- устоявшаяся исследовательская задача с~собственными датасетами и~benchmark'ами:
\begin{itemize}
\item \textbf{AVA (Aesthetic Visual Analysis)}~\cite{Murray:2012} "--- 250\,000 фотографий с~эстетическими оценками; используется для~aesthetic scoring, не~для~стилевой классификации.
\item \textbf{Flickr"~Style}~\cite{Karayev:2014} "--- 80\,000 фотографий, размеченных по~20~стилям (HDR, Vintage, Bokeh, Macro и~др.); первая публично доступная модель стилевой классификации на~CNN.
\item Работы по~photographic style transfer~\cite{Luan:2017}, содержательно обсуждающие операциональное определение <<стиля>>.
\end{itemize}

Разрабатываемая система методологически следует Flickr"~Style: дискретная таксономия + supervised classification, но~с~собственной сокращённой и~операционально определённой таксономией восьми классов, что повышает per"=class качество за~счёт меньшего перекрытия категорий.

\paragraph{Позиционирование.} Разрабатываемая система занимает не~занятую коммерческими продуктами нишу: \textbf{open"=source, self"=hosted, специализированный на~стиле классификатор фотоархива с~публичным embedding"=API}. Целевая аудитория "--- фотографы, дизайнеры, кураторы коллекций, исследователи в~области компьютерного зрения, образовательные программы по~глубокому обучению.

\section{Обзор технологий программного стека}
\label{sec:stack}

Каждый компонент технологического стека выбран после явного сравнения с~альтернативами по~критериям зрелости экосистемы, эксплуатационных характеристик и~пригодности для~поставленной задачи. Ниже приведены обоснования по~ключевым компонентам.

\subsection{Backend-фреймворк}

В~качестве кандидатов рассмотрены Spring~Boot~3 (Java~17), Django REST Framework (Python), Express/Fastify (Node.js), Gin/Echo (Go). Выбран \textbf{Spring~Boot~3}.

Главный аргумент "--- \emph{отделение от~Python"=стека ML"=сервиса}. Backend и~ML"=сервис работают на~разных платформах с~независимыми зависимостями, что~исключает конфликты библиотек и~упрощает независимое масштабирование. Зрелая инфраструктура Spring Security (фильтр"=чейн, JWT, BCrypt) <<из~коробки>> снижает риск ошибок в~критичной auth"=логике. Существенным дополнительным фактором является развитая тестовая инфраструктура с~Testcontainers (PostgreSQL, RabbitMQ, MinIO стартуют в~Docker как реальные сервисы для~интеграционных тестов).

\subsection{ML-инференс}

\paragraph{Язык и веб-фреймворк.} Для ML"=сервиса выбран \textbf{FastAPI} на~Python~3.11. Альтернативы "--- Flask (синхронный по~умолчанию), TorchServe и~Triton Inference Server (production"=решения, оптимизированные для~высоконагруженных сервисов на~GPU). Для~дипломной работы важна прозрачность инференс"=кода (студент должен видеть и~контролировать каждый шаг) при~минимальной кодовой базе; FastAPI обеспечивает все необходимые возможности (HTTP"~API + фоновые consumer'ы AMQP) и~автогенерирует OpenAPI"=документацию.

\paragraph{Фреймворк глубокого обучения.} Выбран \textbf{PyTorch~2.4}. В~research"=публикациях после~2020 года PyTorch является стандартом de~facto: более 80\,\% статей CVPR 2023 используют именно его. Это означает максимально широкий выбор предобученных моделей и~учебных материалов. Библиотека \texttt{torchvision.models} напрямую отдаёт EfficientNet"~B0 с~ImageNet"=весами одной строкой.

\subsection{Очередь сообщений}

В~качестве кандидатов рассмотрены RabbitMQ, Apache Kafka, Redis Streams, Amazon SQS, ZeroMQ. Выбран \textbf{RabbitMQ}.

AMQP"=протокол предоставляет ровно нужный набор примитивов: durable очереди с~гарантированной доставкой, direct exchange с~routing keys для~маршрутизации, acknowledgements для~подтверждения обработки. Management"=UI существенно облегчает отладку. Apache Kafka переинженерен для~сценария одного producer'а и~одного consumer'а с~низким TPS "--- его архитектурные преимущества (партиционирование, replay'и) в~настоящей задаче не~востребованы. Redis Streams "--- молодая технология с~ограниченным набором функций маршрутизации.

\subsection{Объектное хранилище}

В~качестве кандидатов рассмотрены MinIO, AWS~S3, локальная файловая система, Cloudinary. Выбран \textbf{MinIO}.

Главный аргумент "--- \emph{API"=совместимость с~AWS~S3}: тот~же \texttt{boto3} (Python) или~\texttt{minio"=java} SDK работает с~обоими хранилищами; миграция на~AWS~S3 при~необходимости production"=развёртывания заключается в~смене endpoint'а в~конфигурации, без~изменений в~коде. Локальная файловая система отвергнута из"~за отсутствия HA"=гарантий и~непригодности для~горизонтально масштабируемого backend'а. Cloudinary "--- закрытый коммерческий сервис, не~допускающий self"=hosting.

\subsection{Реляционная база данных}

В~качестве кандидатов рассмотрены PostgreSQL~16, MySQL/MariaDB, MongoDB, SQLite. Выбран \textbf{PostgreSQL}.

Сильная типизация защищает от~порчи данных, ACID"=гарантии важны для~поддержания согласованности тройки \texttt{photos $\leftrightarrow$ photo\_styles $\leftrightarrow$ photo\_features}. Дополнительный фактор: для~планируемого масштабирования embedding"=поиска возможен переход на~расширение \texttt{pgvector}, что~позволит сохранить хранилище и~миграции в~рамках PostgreSQL.

\subsection{Клиентский фреймворк}

В~качестве кандидатов рассмотрены React~18, Vue~3, Angular, Svelte. Выбран \textbf{React~18 + Vite~5}.

React не~является лучшим по~техническим характеристикам в~абсолютном выражении (Svelte и~Vue~3 имеют меньший runtime, более полную типизацию), но~является самой востребованной технологией на~рынке труда "--- что~является существенным образовательным аргументом для~дипломной работы. Сборщик Vite обеспечивает Hot Module Replacement и~быстрый dev"=сервер, превышающий по~производительности Webpack/CRA в~5"=10 раз.

\subsection{Аутентификация}

В~качестве кандидатов рассмотрены JWT (stateless), session"=cookies со~стороне"=сервера, делегирование OAuth2 третьей стороне. Выбран \textbf{JWT~(HS512, TTL~24~ч.)}.

Stateless"=архитектура backend'а делает возможным горизонтальное масштабирование без~необходимости shared session store. Отзыв токенов в~текущей версии не~реализован (приемлемо для~прототипа); в~production"=развёртывании для~отзыва добавляется blacklist в~Redis или~короткие TTL с~refresh"=токенами.

\subsection{Оркестрация развёртывания}

Выбран \textbf{Docker~Compose}. Для~дипломной демонстрации Kubernetes избыточен; одной командой \texttt{docker compose up} поднимаются пять контейнеров с~корректными зависимостями старта и~общей сетью. Перенос на~Kubernetes "--- добавление пяти манифестов Deployment+Service на~основе существующих Dockerfile'ов; образы переиспользуются.

%%%%%%%%%%%%%%%%%%%%%%%%%%%%%%%%%%%%%%%%%%
%% ГЛАВА 2
%%%%%%%%%%%%%%%%%%%%%%%%%%%%%%%%%%%%%%%%%%
\chapter{Проектирование системы}
\label{ch:design}

\section{Архитектура системы}
\label{sec:architecture}

Система спроектирована по~микросервисному принципу. Выделено пять независимых компонентов с~чёткими границами ответственности (табл.~\ref{tab:components}). Взаимодействие выполняется по~двум типам каналов: синхронному (HTTP+JSON между клиентом и~backend'ом, между backend'ом и~хранилищами) и~асинхронному (AMQP"=сообщения между backend'ом и~ML"=сервисом через брокер RabbitMQ).

\begin{table}[!ht]
\caption{Компоненты системы и их назначение}
\label{tab:components}
\centering
\begin{tabular}{p{3.5cm}p{2cm}p{8cm}}
\toprule
\textbf{Компонент} & \textbf{Платформа} & \textbf{Назначение} \\
\midrule
React SPA            & Vite + браузер   & Регистрация, загрузка фотографий, отображение галереи, модальное окно с~анализом стилей и~похожими фотографиями \\
Spring Backend       & JVM (Java~17)    & REST"=API, авторизация (JWT), управление учётными записями, оркестрация задач классификации, поиск по~стилю и~поиск похожих изображений по~эмбеддингу \\
FastAPI ML"=сервис   & Python~3.11      & Загрузка обученной модели, инференс на~изображениях из~MinIO, извлечение эмбеддингов и~палитры, публикация результатов в~очередь \\
PostgreSQL~16        & --                & Хранение пользователей, метаданных фотографий, связей фотография"=стиль с~уверенностями, эмбеддингов и~score"=card"=метрик \\
MinIO                & --                & Объектное хранилище загруженных изображений; backend публикует presigned"=URL"ы для~отображения в~клиенте \\
RabbitMQ             & --                & Брокер для~асинхронной передачи задач классификации и~результатов между Spring и~FastAPI \\
\bottomrule
\end{tabular}
\end{table}

Концептуально архитектура реализует трёхуровневое разделение: уровень представления (клиент), уровень приложения (backend и~ML"=сервис) и~уровень хранения (PostgreSQL, MinIO, RabbitMQ как stateful"=компоненты инфраструктуры). Все пять компонентов запускаются в~отдельных контейнерах Docker; их~жизненный цикл управляется через Docker Compose с~зависимостями старта и~healthcheck'ами.

\section{Поток данных при классификации фотографии}
\label{sec:dataflow}

Полный жизненный цикл одной загруженной фотографии включает следующие шаги:

\begin{enumerate}[label=\arabic{enumi})]
\item Пользователь загружает файл изображения через React"=форму запросом \verb|POST /api/photos| с~Content"~Type \verb|multipart/form-data| и~заголовком \verb|Authorization: Bearer <JWT>|.
\item Spring backend выполняет валидацию: проверка подписи JWT, проверка MIME"=типа (whitelist \verb|image/jpeg|, \verb|image/png|, \verb|image/webp|), проверка размера (лимит~25~МБ через \texttt{spring.servlet.multipart}).
\item Файл сохраняется в~MinIO под ключом \verb|user-{userId}/{uuid}.{ext}|. Префикс с~идентификатором пользователя предотвращает горизонтальную утечку privileges, UUID гарантирует уникальность.
\item Создаётся запись в~таблице \texttt{photos} со~статусом \texttt{PENDING}.
\item Backend публикует сообщение в~exchange \verb|classification.exchange| с~routing key \verb|classification.task|, содержащее $\{photoId, s3Key, bucket\}$.
\item ML"=сервис получает сообщение из~очереди \verb|classification.tasks|, скачивает файл из~MinIO по~$(bucket, s3Key)$, выполняет препроцессинг (resize до~256, центрированный crop до~$224\times 224$, нормализация ImageNet"=статистиками) и~прямой проход через сеть.
\item Выход сети разделяется на~два потока: top"~3 стилевых ярлыка с~постмягкомаксными вероятностями и~1280"=мерный embedding с~предклассифицирующего слоя.
\item Результат, обогащённый палитрой (5 цветов через k"~means на~уменьшенной копии) и~score"~card"=метриками (5~детерминированных характеристик), публикуется в~очередь \verb|classification.results|.
\item Spring backend (consumer) получает результат и~в~рамках одной транзакции: удаляет прежние записи \texttt{photo\_styles} по~$photoId$, вставляет новые, сохраняет embedding в~\texttt{photo\_features}, переводит \verb|photos.status| в~\texttt{DONE}.
\item Клиент, выполняющий периодический опрос \verb|GET /api/photos| каждые три~секунды при~наличии \texttt{PENDING}"=фотографий, получает обновлённое состояние и~перерисовывает галерею.
\end{enumerate}

Полный цикл от~момента нажатия кнопки загрузки до~отображения тегов в~галерее занимает порядка $2$"=$3$ секунд для~файла размером 2~МБ при~CPU"=инференсе.

\section{Модель данных}
\label{sec:datamodel}

Реляционная схема включает пять основных сущностей. Их~отношения отражены в~листинге~\ref{lst:schema}.

\begin{lstlisting}[caption={Схема реляционной модели данных},label=lst:schema,basicstyle=\ttfamily\small,frame=single]
users        (id PK, email UNIQUE, password_hash, created_at)
styles       (id PK, name UNIQUE)
photos       (id PK, user_id FK->users, s3_key UNIQUE,
              uploaded_at, status ENUM('PENDING','DONE','FAILED'))
photo_styles (photo_id FK->photos, style_id FK->styles,
              confidence DECIMAL(5,4),
              PRIMARY KEY (photo_id, style_id))
photo_features (photo_id PK/FK->photos, embedding TEXT,
                palette JSON, scores JSON, brightness REAL,
                contrast REAL, saturation REAL,
                warmth REAL, sharpness REAL,
                dominant_hex TEXT)
\end{lstlisting}

Жизненный цикл фотографии описывается конечным автоматом с~тремя состояниями: \texttt{PENDING} (после загрузки, в~очереди), \texttt{DONE} (успешная классификация), \texttt{FAILED} (ошибка в~ML"=сервисе либо таймаут). Переходы возможны только из~\texttt{PENDING} в~два других состояния; обратные и~боковые переходы не~допускаются.

Справочник \texttt{styles} инициализируется при~старте через миграции Flyway (\texttt{V1\_\_init.sql}, \texttt{V2}, ..., \texttt{V6\_\_taxonomy\_v3\_8\_classes.sql}). Хранение в~таблице, а~не~в~enum, выбрано для~поддержки нескольких итераций таксономии без~изменения схемы Java"=кода.

Поле \verb|photo_features.embedding| хранит 1280"=мерный вектор в~виде JSON"=строки. При~планируемом масштабировании за~$\sim$\,100\,000~изображений предусмотрен переход на~расширение \texttt{pgvector}, при~котором тип столбца меняется на~\texttt{vector(1280)}, а~поиск ближайших соседей выполняется через индекс HNSW. Текущая реализация выполняет поиск брутфорсно в~Java"=коде сервиса \texttt{SimilarityService} (cosine similarity, $O(N \cdot 1280)$); для~коллекций до~$\sim$\,10\,000~изображений время отклика остаётся менее~50~мс.

\section{REST API}
\label{sec:api}

API спроектирован по~принципам REST с~JSON"=форматом сериализации. Эндпоинты сгруппированы по~четырём смысловым семействам: аутентификация, управление фотографиями, поиск и~справочные данные (табл.~\ref{tab:endpoints}).

\begin{table}[!ht]
\caption{Эндпоинты REST"-API}
\label{tab:endpoints}
\centering
\begin{tabular}{lp{5cm}p{4cm}c}
\toprule
\textbf{Метод} & \textbf{Путь} & \textbf{Описание} & \textbf{JWT} \\
\midrule
POST   & \texttt{/api/auth/register}     & Регистрация пользователя   & нет \\
POST   & \texttt{/api/auth/login}        & Аутентификация             & нет \\
POST   & \texttt{/api/photos}            & Загрузка фото (multipart)  & да  \\
GET    & \texttt{/api/photos}            & Список фото с~пагинацией   & да  \\
GET    & \texttt{/api/photos/\{id\}}     & Одна фотография            & да  \\
DELETE & \texttt{/api/photos/\{id\}}     & Удаление                   & да  \\
GET    & \texttt{/api/photos/\{id\}/similar} & Похожие по~embedding   & да  \\
GET    & \texttt{/api/photos/search}     & Поиск по~стилю/цвету       & да  \\
GET    & \texttt{/api/styles}            & Список стилей              & да  \\
GET    & \texttt{/api/stats/me}          & Статистика пользователя    & да  \\
\bottomrule
\end{tabular}
\end{table}

Все ответы возвращаются в~JSON. Унифицированная структура ошибок реализована через \texttt{ApiError}, содержащий поля \verb|timestamp|, \verb|status|, \verb|error|, \verb|message|, \verb|path|. Глобальный обработчик \texttt{GlobalExceptionHandler} перехватывает типовые исключения и~отображает их~в~HTTP"=статусы: \texttt{MethodArgumentNotValidException}~$\to$~400, \texttt{BadCredentialsException}~$\to$~401, \texttt{IllegalStateException}~$\to$~409, \texttt{IllegalArgumentException}~$\to$~400. Доступ без~корректного JWT отдаёт~401 (через явно настроенный \texttt{HttpStatusEntryPoint}).

\section{Аутентификация и авторизация}
\label{sec:auth}

Применена JWT"=аутентификация без~серверных сессий. Пароли хранятся в~виде BCrypt"=хешей со~strength~10. При~успешной аутентификации сервер генерирует токен с~алгоритмом подписи HS512, сроком жизни 24~часа, содержащий \texttt{sub}~=~email пользователя. Секретный ключ передаётся через переменную окружения \verb|JWT_SECRET|; в~репозитории \texttt{application.yml} содержит заведомо невалидное значение \verb|change-me-in-prod|.

Каждый последующий запрос несёт заголовок \verb|Authorization: Bearer <token>|. Фильтр \texttt{JwtAuthFilter}, размещённый в~цепочке перед \texttt{UsernamePasswordAuthenticationFilter}, извлекает токен, проверяет подпись и~срок действия, помещает \texttt{UserDetails} в~\texttt{SecurityContext}. Все запросы к~ресурсам пользователя фильтруются по~\verb|user_id| на~уровне SQL"=запроса в~репозитории (\verb|findByIdAndUserId|), что~исключает horizontal privilege escalation на~уровне кода, а~не~только бизнес"=логики.

%%%%%%%%%%%%%%%%%%%%%%%%%%%%%%%%%%%%%%%%%%
%% ГЛАВА 3
%%%%%%%%%%%%%%%%%%%%%%%%%%%%%%%%%%%%%%%%%%
\chapter{Реализация}
\label{ch:implementation}

\section{Структура проекта}
\label{sec:structure}

Репозиторий организован в~корневые каталоги по~компонентам системы (листинг~\ref{lst:tree}). Каждый компонент содержит собственный Dockerfile и~конфигурацию сборки. Корневой \verb|docker-compose.yml| описывает их~совместный запуск.

\begin{lstlisting}[caption={Структура корневого каталога репозитория},label=lst:tree,basicstyle=\ttfamily\small,frame=single]
photo-style-classifier/
|-- docker-compose.yml
|-- README.md
|-- backend/             Java/Spring Boot 3
|   |-- Dockerfile
|   |-- pom.xml
|   `-- src/main/java/com/diploma/psc/
|       |-- auth/         JWT, контроллер, сервис, фильтр
|       |-- user/         User + репозиторий
|       |-- photo/        Photo, MinioService, SimilarityService
|       |-- style/        Style, PhotoStyle, репозитории
|       |-- classification/ Producer, Consumer, DTO
|       |-- stats/        StatsController, StatsResponse
|       |-- config/       Security, Minio, RabbitMQ
|       `-- common/       GlobalExceptionHandler
|-- ml-service/          Python/FastAPI
|   |-- Dockerfile
|   |-- requirements.txt
|   |-- weights/         .pth с обученной моделью
|   `-- app/
|       |-- config.py    pydantic-settings
|       |-- classifier.py StyleClassifier, EfficientNet-B0
|       |-- minio_client.py boto3-загрузчик
|       |-- rabbitmq_consumer.py pika-worker
|       `-- main.py      FastAPI app + handle_task
|-- frontend/            React + Vite
|   |-- package.json
|   |-- vite.config.js
|   `-- src/
|       |-- App.jsx, main.jsx, AuthContext.jsx, api.js
|       |-- routes/       Login, Register, Gallery, Upload, Profile
|       `-- components/   PhotoCell, PhotoModal, ColorFilterBar
|-- training/            Скрипты обучения и экспериментов
|-- docs/                REPORT.md, diagrams.md, training-metrics.md
`-- .github/workflows/   CI pipelines
\end{lstlisting}

\section{Backend на Spring Boot}
\label{sec:backend}

\subsection{Управление схемой через Flyway и режим JPA validate}

Схема создаётся и~изменяется только миграциями Flyway, хранимыми в~\verb|backend/src/main/resources/db/migration/|. Hibernate настроен в~режиме \verb|ddl-auto: validate|: при~старте проверяется соответствие сущностей таблицам и~приложение падает при~рассогласовании. Это исключает ситуации тихого изменения production"=схемы при~деплое и~делает миграции единственным источником правды о~структуре БД.

\subsection{Двойной клиент MinIO}

В~конфигурации определены два бина \texttt{MinioClient} с~разными endpoint'ами. Основной \texttt{minioClient} использует адрес \verb|http://minio:9000| (имя сервиса в~Docker"=сети) и~применяется для~операций \verb|putObject|, \verb|bucketExists|, \verb|makeBucket|, выполняемых внутри Docker"=сети. Дополнительный \texttt{minioPresignClient} использует адрес \verb|http://localhost:9000| с~фиксированным регионом \verb|us-east-1| и~применяется только для~генерации presigned"=URL'ов, отдаваемых клиенту: браузер не~способен разрешить имя \verb|minio:9000|, поэтому URL должен содержать хост, доступный с~host"=машины.

Фиксация региона в~presign"=клиенте критически важна: без~неё SDK при~первом вызове выполняет HTTP"=запрос \verb|GET ?location| к~бакету, что~приводит к~попытке обращения к~\verb|localhost:9000| из~контейнера и~ошибке.

\subsection{Асинхронная публикация через RabbitTemplate}

Exchange типа \verb|direct| с~routing key для~маршрутизации задач классификации. Сообщения сериализуются через \texttt{Jackson2JsonMessageConverter}, что~позволяет Python"=консьюмеру работать с~ними как с~обычным JSON без~необходимости поддержки сериализаторов Java.

Consumer на~стороне backend'а получает результаты классификации и~выполняет атомарное обновление: сначала удаляются все строки \texttt{photo\_styles} по~$photoId$ (вызов \verb|flush()| для~гарантии порядка), затем вставляются новые. Это защищает от~некорректных состояний при~повторной обработке (retry), когда часть старых тегов могла бы остаться в~БД.

\subsection{Обработка ошибок}

\texttt{GlobalExceptionHandler} с~аннотацией \verb|@RestControllerAdvice| унифицирует ответы при~неуспешных запросах. Это позволяет фронтенду полагаться на~единый формат и~централизовать обработку ошибок в~одном axios"=interceptor'е.

\section{ML-сервис на FastAPI}
\label{sec:mlservice}

\subsection{Построение модели}

Архитектура модели создаётся вызовом \verb|efficientnet_b0(weights=None)| с~последующей заменой классифицирующего слоя:

\begin{lstlisting}[language=Python,basicstyle=\ttfamily\small,frame=single,caption={Создание модели и подключение весов},label=lst:model]
from torchvision.models import efficientnet_b0
model = efficientnet_b0(weights=None)
in_features = model.classifier[1].in_features  # 1280
model.classifier[1] = nn.Linear(in_features, len(self.styles))
\end{lstlisting}

При~наличии файла \verb|weights/efficientnet_b0_styles.pth| в~него загружается \texttt{state\_dict} дообученной модели. Файлы весов исключены из~репозитория (\verb|.gitignore|): они являются крупными бинарниками с~независимым жизненным циклом и~хранятся в~Google~Drive автора либо передаются вручную при~деплое.

При~отсутствии файла весов сервис стартует в~\emph{эвристическом режиме}: вычисляет статистики изображения (средняя яркость, контраст, насыщенность, теплота, динамический диапазон) и~маппит их~в~вероятности классов через набор взвешенных правил. Этот режим предназначен для~end"=to"=end отладки всей системы при~отсутствии обученной модели и~не~является заменой обученной модели.

\subsection{Извлечение эмбеддингов и сплит forward-прохода}

Ключевая особенность реализации "--- расщепление прямого прохода на~два этапа для~извлечения эмбеддингов вместе со~стилевыми вероятностями:

\begin{lstlisting}[language=Python,basicstyle=\ttfamily\small,frame=single,caption={Расщеплённый forward-проход},label=lst:fwd]
with torch.no_grad():
    x = self.transform(img).unsqueeze(0).to(self.device)
    feats = self.model.features(x)
    pooled = self.model.avgpool(feats)
    pooled_flat = torch.flatten(pooled, 1)        # [1, 1280] - эмбеддинг
    logits = self.model.classifier(pooled_flat)   # [1, num_classes]
    probs = torch.softmax(logits, dim=1)
    embedding = pooled_flat.squeeze(0).cpu().numpy()
\end{lstlisting}

Эмбеддинг возвращается из~предклассифицирующего слоя (выход global average pooling), классифицирующая голова применяется отдельно. Это позволяет получить оба выхода за~один проход и~опубликовать эмбеддинг вместе с~результатом классификации, не~требуя повторного forward'а.

\subsection{Препроцессинг и аугментации}

Препроцессинг детерминированный, без~аугментаций: resize до~$224\times 224$ (точно, без~CenterCrop, чтобы сохранить пропорции), приведение к~\texttt{torch.float32} в~диапазоне $[0;1]$, нормализация ImageNet"=статистиками ($\mu = [0{,}485; 0{,}456; 0{,}406]$, $\sigma = [0{,}229; 0{,}224; 0{,}225]$). Идентичный препроцессинг применяется и~при~обучении (см.~раздел~\ref{sec:training}), чтобы избежать train/serve skew.

\subsection{Цветовая палитра и score-card}

Помимо стилевой классификации, ML"=сервис извлекает два дополнительных слоя описания изображения. Оба вычисляются детерминированно из~пиксельных статистик, без~участия нейросети.

\paragraph{Палитра цветов} "--- 5 доминирующих цветов в~RGB через \texttt{sklearn.cluster.KMeans} (\verb|n_clusters=5|, \verb|n_init=3|) на~уменьшенной до~$100\times 100$ копии изображения. Кластеры сортируются по~размеру (наиболее частый цвет первым) и~возвращаются как hex"=строки. Доминирующий hex"=цвет дополнительно сохраняется в~\verb|photo_features.dominant_hex| для~быстрой фильтрации галереи по~цвету.

\paragraph{Score-card} "--- 5 количественных метрик в~диапазоне $[0;1]$:

\begin{itemize}
\item \texttt{brightness} "--- средняя perceptual luma $Y' = 0{,}299 R + 0{,}587 G + 0{,}114 B$;
\item \texttt{contrast} "--- стандартное отклонение luma, удвоенное и~ограниченное на~$1{,}0$;
\item \texttt{saturation} "--- среднее значение нормированной формулы $(max - min) / max$ в~RGB;
\item \texttt{warmth} "--- нормализованная разность $R - B$ в~диапазон $[0;1]$;
\item \texttt{sharpness} "--- средняя абсолютная величина градиента luma, как proxy для~резкости.
\end{itemize}

Эти метрики играют двойную роль. На~уровне UI они дополняют стилевые ярлыки описательной информацией. На~уровне методологии они позволяют формально определить классы через измеримые статистики (см.~раздел~\ref{sec:scorecard}) и~использовать как baseline"=признаки для~оценки вклада глубоких представлений.

\subsection{AMQP-консьюмер в фоновом потоке}

Библиотека \texttt{pika} использует блокирующее соединение, что~несовместимо с~asyncio. Решение "--- запуск consumer'а в~отдельном \texttt{threading.Thread} с~обработкой \texttt{AMQPConnectionError} (пауза 5~с, новая попытка). FastAPI"=приложение стартует независимо и~предоставляет endpoint \verb|/health| для~диагностики (возвращает имя модели, путь к~весам, имя устройства, список стилей).

Префетч очереди ограничен двумя задачами через \verb|channel.basic_qos(prefetch_count=2)|: инференс одного изображения требует $\sim$\,500~МБ peak RSS, и~накопление задач выше двух привело~бы к~OOM при~всплеске загрузок.

\section{Клиентское приложение на React}
\label{sec:frontend}

\subsection{Дизайн}

Применён тёмный brutalist"=стиль галереи: фон \verb|#111|, поверхности \verb|#141414|, акцентный цвет \verb|#c0392b|, отсутствие скруглений, типографика \emph{uppercase} с~расширенным трекингом, мелкий кегль. Цель оформления "--- не~отвлекать от~содержимого (фотографий) и~подчеркнуть утилитарный характер инструмента.

\subsection{Masonry-сетка}

Галерея организована в~детерминированную masonry"=сетку через CSS~Grid: \verb|grid-template-columns: repeat(6, 1fr)|, \verb|grid-auto-rows: 180px|, \verb|grid-auto-flow: dense|. Каждой ячейке на~основе её~индекса присваивается класс из~детерминированного паттерна (\verb|w2|, \verb|h2|, \verb|w2 h2|), что~создаёт визуально неоднородный ритм без~использования случайности. Опция \verb|dense| гарантирует, что~крупные ячейки не~оставляют пустот.

\subsection{Реактивный опрос статуса}

После загрузки фотография имеет статус \texttt{PENDING}. Клиент в~\verb|useEffect| выполняет повторный запрос \verb|GET /api/photos| каждые три~секунды, пока в~ответе присутствует хотя~бы одна \texttt{PENDING}"=фотография. Это простой trade"=off против WebSocket/SSE: нагрузка на~сервер минимальна, код "--- несколько строк в~React"=компоненте, не~требует дополнительной инфраструктуры на~бэкенде.

\subsection{Axios-interceptor для централизованной обработки 401/403}

При~получении кодов 401 или~403 от~любого endpoint'а взаимодействия токен очищается из~\verb|localStorage|, и~пользователь перенаправляется на~маршрут \verb|/login|. Это унифицирует обработку истечения JWT без~дублирования логики в~каждом компоненте, обращающемся к~API.

\subsection{Модальное окно с похожими фотографиями}

При~клике на~ячейку галереи открывается модальное окно с~увеличенной фотографией, цветовой палитрой, score"=card"=метриками, top"~3 стилевыми ярлыками с~прогресс"=барами и~сеткой из~шести похожих фотографий, найденных через \verb|GET /api/photos/{id}/similar|. Похожесть определяется на~основе cosine similarity между 1280"=мерными эмбеддингами (см.~раздел~\ref{sec:embed_quality}). Клик по~похожей фотографии перенаправляет на~неё, что~позволяет навигировать по~коллекции через стилевое <<сходство>>.

\section{Инфраструктура развёртывания}
\label{sec:infrastructure}

\verb|docker-compose.yml| описывает шесть сервисов: \texttt{postgres}, \texttt{rabbitmq}, \texttt{minio}, \texttt{minio-init} (одноразовый инициализатор бакета при~первом запуске), \texttt{backend}, \texttt{ml-service}. Healthcheck'и и~зависимости через \verb|depends_on: condition: service_healthy| гарантируют корректный порядок старта. Все сервисы помещены в~общую Docker"=bridge"=сеть \texttt{psc-network}, что~обеспечивает разрешение имён хостов внутри.

Именованные тома \texttt{postgres\_data}, \texttt{rabbitmq\_data}, \texttt{minio\_data} обеспечивают сохранность данных между перезапусками контейнеров. Frontend в~production"=сборке мог~бы быть упакован в~отдельный nginx"=контейнер; в~настоящей дипломной демонстрации он~запускается в~dev"=режиме Vite через \verb|npm run dev|.

%%%%%%%%%%%%%%%%%%%%%%%%%%%%%%%%%%%%%%%%%%
%% ГЛАВА 4
%%%%%%%%%%%%%%%%%%%%%%%%%%%%%%%%%%%%%%%%%%
\chapter{Датасет и обучение модели}
\label{ch:dataset}

\section{Сбор датасета}
\label{sec:collection}

Источниками изображений выбраны общедоступные стоковые платформы Unsplash и~Pexels, предоставляющие бесплатное API с~достаточным лимитом запросов и~свободной лицензией использования. Эти источники характеризуются высоким техническим качеством снимков и~достаточно широким спектром стилей, что~критично для~стилевой классификации.

Скачивание выполнено через Python"=скрипт \verb|01_dataset_collection.ipynb|, обращающийся к~Unsplash Search API и~Pexels Search API. Для~каждого целевого класса задавался набор поисковых запросов на~английском языке (например, для~\texttt{neon}: <<neon photography>>, <<cyberpunk lights>>, <<synthwave aesthetic>>, <<tokyo neon>> и~др.). На~каждый запрос API возвращает страницу с~30~изображениями; страницы перебираются до~достижения целевого количества $\sim$\,250 фотографий с~каждого источника на~класс.

Изображения сохраняются с~префиксом источника (\verb|u_| для~Unsplash, \verb|p_| для~Pexels) и~уникальным идентификатором: это упрощает диагностику и~отслеживание происхождения. На~первой итерации собрано $\sim$\,1000~изображений на~класс; распределение по~классам оказалось неравномерным (от~491 в~\texttt{monochrome} до~1457 в~\texttt{minimalist}) в~силу различной популярности тематик в~стоковых архивах.

\section{Эволюция таксономии}
\label{sec:taxonomy}

Финальная таксономия восьми классов "--- результат трёх итераций пересмотра, мотивированных диагностикой обученных моделей. История изменений приведена в~таблице~\ref{tab:taxonomy}.

\begin{table}[!ht]
\caption{Эволюция таксономии и сводные метрики}
\label{tab:taxonomy}
\centering
\begin{tabular}{lcccc}
\toprule
\textbf{Итерация} & \textbf{Классы} & \textbf{Best val acc} & \textbf{max/min} & \textbf{Bias} \\
\midrule
v1 (наивный fine-tune) & 8 (с~\texttt{moody}, \texttt{street}) & 0{,}678 & $5{,}0\times$ & сильный на~\texttt{minimalist} \\
v2 (балансировка)      & 7 (без~\texttt{street})              & 0{,}724 & $1{,}9\times$ & заметный на~\texttt{moody} \\
v3 (новая таксономия)  & 8 ($+$\texttt{monochrome}, \texttt{neon}, $-$\texttt{moody}) & 0{,}726 & $1{,}2\times$ & устранён \\
\textbf{v3.1 (воспр.)} & \textbf{та~же}                       & \textbf{0{,}694} & \textbf{$1{,}2\times$} & \textbf{устранён} \\
\bottomrule
\end{tabular}
\end{table}

\paragraph{v1.} Первоначальный набор из~восьми классов: \texttt{moody}, \texttt{minimalist}, \texttt{street}, \texttt{golden\_hour}, \texttt{dark}, \texttt{airy}, \texttt{vintage}, \texttt{dramatic}. Обучение без~балансировки на~датасете с~max/min$=5{,}0$ дало val~acc$=0{,}678$ при~систематическом смещении предсказаний в~доминирующий класс \texttt{minimalist} (ratio предсказанных к~истинным $\sim$\,$2{,}5$).

\paragraph{v2.} Удалён класс \texttt{street} как жанровая категория, не~являющаяся визуальным стилем (модель училась shortcut'у: <<люди на~улице>>~$\to$~\texttt{street}). Применена балансировка через WeightedRandomSampler. Val~acc выросла до~$0{,}724$, но~остаточный bias на~класс \texttt{moody} сохранился. Анализ confusion matrix показал, что~\texttt{moody} систематически пересекается с~\texttt{dark}, \texttt{dramatic}, \texttt{vintage} "--- то~есть является не~ортогональным, а~композитным классом без~операционального определения.

\paragraph{v3.} Класс \texttt{moody} удалён. Введены два класса с~явными операциональными определениями: \texttt{monochrome} (saturation~$\approx 0$) и~\texttt{neon} (saturation~$\geq 0{,}6$ при~низкой средней яркости). Применена двухфазная схема обучения и~MD5"=дедупликация. Val~acc=$0{,}726$, bias"=эффект устранён: распределение предсказаний по~классам в~коридоре $0{,}87$"=$1{,}08$ от~истинного. Per"=class F1: \texttt{neon}~$=0{,}926$, \texttt{golden\_hour}~$=0{,}856$ как лучшие; \texttt{monochrome}~$=0{,}520$ как слабый из"~за ортогонального пересечения с~композиционно"=определимыми классами (\texttt{minimalist}, \texttt{vintage}, \texttt{dark}).

\paragraph{v3.1.} Воспроизводимое переобучение на~детерминированном split'е. Контекст и~необходимость "--- в~разделе~\ref{sec:reproduce}.

\section{Предобработка и балансировка}
\label{sec:preprocess}

Собранный датасет подвергается трём этапам предобработки:

\begin{enumerate}[label=\arabic{enumi})]
\item \textbf{Cap по~количеству.} Каждый класс случайно усекается до~600 фотографий (с~детерминированным \texttt{random.seed(42)}). Класс \texttt{minimalist} как изначально наиболее раздутый (1457 фото) дополнительно ограничивается до~500 после разделения на~train/val. Класс \texttt{monochrome} с~исходными 491 фотографиями не~подвергается обрезке.
\item \textbf{MD5"=дедупликация.} В~каждом классе вычисляется MD5"=хеш содержимого каждого файла; дубликаты удаляются с~сохранением одной копии. На~исходном датасете обнаружено и~удалено 102~дубликата (наиболее значимо в~\texttt{vintage} "--- 25~файлов, \texttt{airy} "--- 27, \texttt{dramatic} "--- 17). Дубликаты возникают из"~за того, что~одни и~те~же фотографии могут быть проиндексированы по~разным поисковым запросам.
\item \textbf{Stratified split} в~пропорции 80/20 с~детерминированным seed=42. Cap по~\texttt{minimalist} (500~файлов) применяется внутри split'а, что~даёт 400~train / 100~val для~этого класса.
\end{enumerate}

Итоговое распределение по~классам приведено в~таблице~\ref{tab:dataset}. Соотношение max/min для~train"=множества составляет $1{,}22\times$ (\texttt{neon}~$=$~480, \texttt{minimalist}~$=$~391), что~считается приемлемым уровнем дисбаланса.

\begin{table}[!ht]
\caption{Распределение train/val по классам в v3.1}
\label{tab:dataset}
\centering
\begin{tabular}{lrr}
\toprule
\textbf{Класс} & \textbf{Train} & \textbf{Val} \\
\midrule
airy        & 445 & 112 \\
dark        & 464 & 116 \\
dramatic    & 456 & 115 \\
golden\_hour& 469 & 118 \\
minimalist  & 400 & 100 \\
monochrome  & 392 &  99 \\
neon        & 480 & 120 \\
vintage     & 450 & 113 \\
\midrule
\textbf{Итого}  & \textbf{3556} & \textbf{893} \\
\bottomrule
\end{tabular}
\end{table}

\section{Архитектура дообучения}
\label{sec:training}

Принята двухфазная схема fine"=tuning'а EfficientNet"~B0, инициализированной ImageNet"=весами:

\paragraph{Фаза 1 (head only, 5~эпох).} Все веса backbone замораживаются. Обучается только новая классифицирующая голова \verb|Linear(1280, 8)| при~learning~rate $10^{-3}$. Цель фазы "--- быстро адаптировать классификатор к~распределению целевых классов до~того, как~backbone начнёт изменяться. Это предотвращает ранние большие градиенты, которые при~немедленном разморозке могут <<сбить>> предобученные представления.

\paragraph{Фаза 2 (full fine"=tune, до~25~эпох).} Все веса размораживаются. Применяются раздельные learning rate: backbone $10^{-4}$, голова $10^{-3}$ (асимметрия в~десять раз). Используется расписание CosineAnnealingLR с~$T_{\max}=25$. Применяется label smoothing с~$\varepsilon=0{,}1$, что~стабилизирует обучение и~улучшает калибровку. Активирован WeightedRandomSampler с~весами $1/count_c$, что~эффективно даёт каждому классу одинаковую ожидаемую частоту в~минибатче.

\paragraph{Дополнительные параметры.} Размер минибатча~32, оптимизатор AdamW с~weight\_decay$=0{,}01$, mixed precision (\verb|torch.amp.autocast|), early stopping с~patience$=7$~эпох без~улучшения val\_acc. Аугментации train"=множества: RandomResizedCrop(224, scale=(0.7, 1.0)), RandomHorizontalFlip(p=0.5), ColorJitter(brightness=0.15, contrast=0.15, saturation=0.15, hue=0.05). Val"=множество не~аугментируется.

\section{Воспроизводимость v3 → v3.1}
\label{sec:reproduce}

В~ходе подготовки исследовательской части (см.~гл.~\ref{ch:results}) обнаружена методологическая проблема с~оригинальным split'ом v3: он~хранился в~эфемерной файловой системе среды выполнения и~был утрачен после её~перезапуска. При~любой реконструкции split'а на~основе того~же \texttt{random.seed(42)} порядок обхода Drive (FUSE"=mount) оказался нестабильным относительно \texttt{random.shuffle()}, что~приводило к~неполному совпадению с~оригинальным разбиением.

Эмпирически это проявилось в~систематическом завышении val~accuracy с~оригинальных $0{,}726$ до~$0{,}925$ при~прогоне той~же модели на~<<реплицированном>> val. Per"=class F1 показали аномальные скачки: \texttt{monochrome}~$0{,}545 \to 0{,}876$, \texttt{dramatic}~$0{,}586 \to 0{,}897$ "--- увеличения на~28"=35~пп. Это характерный отпечаток data leakage: фотографии, которые модель видела в~обучении оригинального v3, после нового split'а оказались в~новом val"=множестве; модель их~<<запомнила>> и~даёт почти безошибочные предсказания, искусственно завышая метрику.

\textbf{Решение.} Полное переобучение на~воспроизводимом split'е. Алгоритм split'а формализован в~скрипте \verb|training/make_split.py| с~детерминированным порядком обхода файлов через \verb|sorted()| и~тем~же seed=42. Гиперпараметры обучения идентичны v3. Финальные веса сохранены как v3.1.

\textbf{Результат.} val~acc$=0{,}694$, macro F1$=0{,}687$. Снижение на~3~пп относительно v3 (от~$0{,}726$) лежит в~пределах статистической дисперсии для~val"=множества из~893 фотографий и~стохастичных факторов обучения (RandomResizedCrop, перетасовка минибатчей, случайная инициализация головы).

Поэпохная динамика обучения v3.1 приведена в~таблице~\ref{tab:training}.

\begin{table}[!ht]
\caption{Поэпохная динамика обучения v3.1}
\label{tab:training}
\centering
\begin{tabular}{ccrrr}
\toprule
\textbf{Фаза} & \textbf{Эпоха} & \textbf{Train Loss} & \textbf{Val Loss} & \textbf{Val Acc} \\
\midrule
1 (head)  &  1 & 1{,}619 & 1{,}406 & 0{,}596 \\
1 (head)  &  3 & 1{,}291 & 1{,}292 & 0{,}639 \\
1 (head)  &  5 & 1{,}240 & 1{,}286 & 0{,}622 \\
2 (full)  &  1 & 1{,}144 & 1{,}205 & 0{,}669 \\
2 (full)  &  3 & 0{,}950 & 1{,}186 & 0{,}679 \\
2 (full)  &  5 & 0{,}846 & 1{,}210 & 0{,}689 \\
2 (full)  & 10 & 0{,}690 & 1{,}248 & 0{,}693 \\
\textbf{2 (full)} & \textbf{12} & \textbf{0{,}699} & \textbf{1{,}256} & \textbf{0{,}694} \\
2 (full)  & 14 & 0{,}656 & 1{,}240 & 0{,}694 \\
2 (full)  & 19 & 0{,}634 & 1{,}266 & 0{,}691 \\
\bottomrule
\end{tabular}
\end{table}

Лучшая эпоха "--- 12 фазы~2 (общая эпоха 17 от~начала обучения). Early stopping сработал на~эпохе~19 фазы~2 после семи эпох без~улучшения val~acc. На~этой эпохе сохранён финальный чекпоинт.

\section{Сохранение весов и деплоймент}
\label{sec:deploy}

После обучения \texttt{state\_dict} лучшей эпохи сохраняется в~файл \verb|efficientnet_b0_styles.pth| в~Google~Drive (каталог \verb|training_out_v3.1/|). Размер чекпоинта составляет $\sim$\,16~МБ. Перенос в~ML"=сервис выполняется ручным копированием в~\verb|ml-service/weights/| и~перезапуском контейнера \verb|docker compose restart ml-service|.

При~старте ML"=сервиса \texttt{StyleClassifier} автоматически обнаруживает файл весов через \verb|_resolve_weights_path()|, который последовательно проверяет: каноничный путь из~конфигурации, файлы с~приоритетным префиксом \verb|efficientnet_b0_styles*|, любой \verb|*.pth| в~директории \verb|weights/|. При~обнаружении валидного state\_dict сервис переходит из~эвристического режима в~режим нейросетевого инференса; событие логируется в~stdout: \verb|Model loaded: EfficientNet-B0 (weights from ...)|.

\end{document}
